\chapter{Adapting the Feed-Forward Head}
\label{ch:adapt}

\section{Choosing the trainable module}
\label{sec:method:head}

Splatt3R trains its Gaussian head against frozen encoder and decoder features.
We compare three adaptation choices:

\begin{itemize}\itemsep2pt
\item \textbf{Encoder LoRA}~\cite{hu2022lora} ($r{=}8$, $\alpha{=}16$ on
      \texttt{qkv}/\texttt{proj}/\texttt{fc1}/\texttt{fc2}):
      $-49\%$ reconstruction quality against the released checkpoint, with
      predicted scales exploding by $1821\times$. The failure is immediate,
      not a late-training instability.
\item \textbf{Decoder-only LoRA}: $+0.37$\dB{} at epoch~1, then collapse to
      $-1.61$\dB{} by epoch~5 with a $12\times$ climb in the $99$th-percentile
      scale.
\item \textbf{Head-only} (encoder and decoder bit-frozen, only
      \texttt{gaussian\_dpt} trainable): improves, and is what we keep.
\end{itemize}

The collapse coincides with the exponential Gaussian-scale parameter
$s_n=\exp(\hat{s}_n)$. A small feature shift can produce a large scale change,
while different scale--opacity combinations can fit the same supervision
views. We therefore retain head-only training for the evaluated system.
The shared encoder, decoder, and geometric heads remain frozen; the Gaussian
DPT head uses the upstream objective and Adam~\cite{kingma2015adam}.

\section{Training setup}
\label{sec:adapt:setup}

Adaptation is performed per dataset family. The families differ in sensor,
exposure, motion, and depth availability. Each receives its own head
checkpoint, loaded at start-up.

The objective is the upstream one of \cref{eq:splatt3rloss}, unchanged, so
that the comparison against the released checkpoint isolates the data and the
choice of what to train, not the loss. Only \texttt{gaussian\_dpt} receives
gradients; the encoder and decoder are frozen, and this is verified by
comparing parameter hashes before and after training rather than by trusting
a \texttt{requires\_grad} flag.

The TUM, EuRoC and 7-Scenes heads in the reported paired adaptation study
were trained for $40$ epochs. Their batch sizes are $2$, $4$ and $8$, with
learning rates $10^{-5}$, $1.41\times10^{-5}$ and $2\times10^{-5}$,
respectively. A $15\%$ validation split is drawn within the available family
data. These experiments measure in-family adaptation; they do not establish
transfer to an unseen dataset family or a sequence-disjoint benchmark.

Two of the five families (EuRoC and ETH3D) ship no usable depth for the
training pairs. Rather than exclude them, pairs are supervised with
self-predicted pseudo-depth from the frozen backbone. This is a weaker
supervision signal, and the results in \cref{ch:results} should be read with
that in mind: EuRoC is the family where the refiner contributes least
($+1.00$\dB) and ETH3D the family where it contributes most ($+8.69$\dB),
and we do not have an explanation for that spread
(\cref{sec:neg:attrib}).

\section{Photometric augmentation and the opacity channel}
\label{sec:adapt:aug}

\Cref{sec:method:mech} examines opacity's photometric role. Inputs darkened
independently of target views provide a controlled way to test whether
opacity statistics respond to appearance perturbations. The SH residual
can also change colour, so this experiment does not isolate opacity as the
only compensating channel.

Four augmentation regimes were trained for $40$ epochs each on the same family
and probed with an identical five-pair protocol. The table reports the
$90$th-percentile predicted scale (a proxy for how far the head has moved from
the released checkpoint's geometry), the median predicted opacity, and the
fraction of Gaussians with opacity above $0.9$ --- i.e.\ how saturated the
opacity channel is:

\begin{center}
\small
\begin{tabular}{@{}lrrrr@{}}
\toprule
Head & scale p90 & inflation & median $\alpha$ & $\alpha>0.9$ \\
\midrule
base (released)              & \phantom{0}9.78 & ---       & 1.000 & 100.0\% \\
content noise + blur         & 26.18 & $+167.8\%$ & 1.000 & 100.0\% \\
global photometry            & 17.37 & $+77.6\%$  & 0.991 & \phantom{0}83.1\% \\
spatial photometry           & 14.82 & $+51.5\%$  & 0.876 & \phantom{0}42.9\% \\
content + spatial photometry & 10.97 & $+12.2\%$  & 0.933 & \phantom{0}62.7\% \\
\bottomrule
\end{tabular}
\end{center}

Three things are visible. The released checkpoint's opacity is
\emph{completely saturated} --- median $1.000$, every Gaussian above $0.9$ ---
which provides little initial differentiation between primitives and
motivates testing an external attenuation prior. Content-only
augmentation does not move the channel at all while inflating scale by
$168\%$. Global photometry modestly reduces saturation; spatial photometry
reduces it much more strongly, from $100\%$ to $42.9\%$. These are observed
parameter responses, rather than a complete account of what the head learned.

The combined regime was tested against an explicit prediction: if the two
perturbations engage independently, saturation should fall to $30$--$50\%$
\emph{and} scale inflation should exceed $10\%$. Scale inflation cleared the
bar at $+12.2\%$, barely; saturation landed at $62.7\%$, \emph{above} the
predicted range and above what spatial photometry achieved on its own. Both
numbers moved in the direction of mutual suppression rather than independent
engagement. The prediction is therefore not supported.

Across four scenes the spatial-photometry
head gives a mean LPIPS improvement of $-14.55\%$, with a per-scene spread of
$-6.6\%$ to $-19.2\%$. The gain is scene-dependent.
